\documentclass[11pt, a4paper]{article}

% ----------------------------------------------------------------------
% PAKETE & EINSTELLUNGEN
% ----------------------------------------------------------------------
\usepackage[utf8]{inputenc}
\usepackage[T1]{fontenc}
\usepackage[ngerman]{babel}
\usepackage{amsmath, amssymb, amsfonts, amsthm}
\usepackage{graphicx}
\usepackage{hyperref}
\usepackage{geometry}
\usepackage{cite}
\usepackage{fancyhdr}
\usepackage{braket}
\usepackage{titlesec}

% Layout-Optimierung für Arxiv-Style (dicht, professionell)
\geometry{left=2cm, right=2cm, top=2cm, bottom=2.5cm}
\setlength{\parindent}{0pt}
\setlength{\parskip}{0.6em}

% Theorem-Umgebungen
\newtheorem{theorem}{Theorem}
\newtheorem{postulate}{Postulat}
\newtheorem{definition}{Definition}

% Titel-Metadaten
\title{\textbf{Die algebraische Notwendigkeit der Riemannschen Vermutung durch spektrale Analyse der Hohlraumstrahlung}

\author{
	\textbfT. Hoffbauer} \\
	\textit{thomas.hoffbauer@protonmail.com} \\
	\texttt{Bamberg, Deutschland}
}

\date{12. Februar 2026}

\begin{document}
	
	\maketitle
	
	\begin{abstract}
		\noindent
		\textbf{Zusammenfassung:} Diese Arbeit verknüpft die analytische Zahlentheorie mit der Quantenstatistik idealer Hohlraumstrahler. Wir postulieren, dass die spektrale Zustandsdichte des Vakuums eine multiplikative Feinstruktur aufweist, die durch die nichttrivialen Nullstellen der Riemannschen Zeta-Funktion determiniert ist. 
		Theoretisch fundiert wird dies durch den \textbf{Satz von Hurwitz} über normierte Divisionsalgebren: Wir zeigen, dass die kritische Linie $\text{Re}(s)=1/2$ keine zufällige analytische Eigenschaft ist, sondern eine algebraische Notwendigkeit, um die Symmetrie zwischen Gaußschen und Eisensteinschen Gittern im 8-dimensionalen Raum (Oktonionen/$E_8$) zu erhalten.
		Empirisch wird das \glqq Bamberg-Protokoll\grqq{} vorgestellt – ein strenges, falsifizierbares Verfahren (Frozen Pipeline), um diese Signatur mittels Gumbel-Statistik aus radiometrischen Messdaten zu extrahieren.
	\end{abstract}
	
	% ----------------------------------------------------------------------
	\section{Einleitung: Das Opfer der Struktur}
	% ----------------------------------------------------------------------
	
	Die Geschichte der modernen Physik ist eine Geschichte des Verlusts klassischer Gewissheiten zugunsten höherer Strukturen. Am 16. Oktober 1843 ritzte William Rowan Hamilton die Formel $i^2 = j^2 = k^2 = ijk = -1$ in die Broom Bridge zu Dublin. Mit der Entdeckung der Quaternionen ($\mathbb{H}$) opferte die Mathematik erstmals die \textit{Kommutativität}, um den dreidimensionalen Raum beschreiben zu können. Wenig später entdeckten Graves und Cayley die Oktonionen ($\mathbb{O}$), womit auch die \textit{Assoziativität} fiel.
	
	Adolf Hurwitz bewies 1898 endgültig, dass es nur vier normierte Divisionsalgebren gibt: $\mathbb{R}, \mathbb{C}, \mathbb{H}$ und $\mathbb{O}$. 
	Parallel dazu entwickelte sich die Quantenmechanik, die Hamiltons Nicht-Kommutativität physikalisch im Heisenbergschen Unschärfeprinzip ($xp \neq px$) manifestierte.
	
	Diese Arbeit stellt die These auf, dass die Riemannsche Vermutung (RH) der direkte algebraische Schatten dieser Hierarchie ist. Wir argumentieren, dass die Nullstellen der Zeta-Funktion die \glqq Nägel\grqq{} sind, die den Übergang zwischen diesen Algebren stabilisieren, und dass diese Stabilität im thermischen Rauschen des Vakuums messbar ist.
	
	% ----------------------------------------------------------------------
	\section{Algebraische Herleitung: Der Hurwitz-Imperativ}
	% ----------------------------------------------------------------------
	
	Bevor wir das Messverfahren beschreiben, leiten wir her, warum die Riemann-Hypothese zwingend wahr sein muss, um die Konsistenz der Divisionsalgebren zu wahren.
	
	\subsection{Gauß, Eisenstein und das $E_8$-Gitter}
	In der komplexen Ebene $\mathbb{C}$ existieren zwei fundamentale Gitterstrukturen für Primzahlen:
	\begin{enumerate}
		\item Das quadratische Gitter der \textbf{Gaußschen Zahlen} $\mathbb{Z}[i]$.
		\item Das hexagonale Gitter der \textbf{Eisenstein-Zahlen} $\mathbb{Z}[\omega]$ ($\omega = e^{i\pi/3}$).
	\end{enumerate}
	In 2D sind diese Symmetrien inkompatibel. Im 8-dimensionalen Raum der Oktonionen jedoch existiert das \textbf{$E_8$-Gitter}, die dichteste Kugelpackung, welche sowohl quaternionische ($D_4$) als auch hexagonale Symmetrien vereinigt.
	Damit eine \glqq Prim-Resonanz\grqq{} (eine Nullstelle $\rho$) universell existiert, muss sie invariant unter dem Basiswechsel zwischen diesen Gittertypen sein. Sie muss \textit{additiv erreichbar} bleiben, wenn man von $\mathbb{H}$ nach $\mathbb{O}$ liftet.
	
	\subsection{Das Hurwitz-Stabilitäts-Funktional}
	Wir definieren die Dualität des Gitters durch die Transformation $s \to 1-s$ (die Funktionalgleichung der Zeta-Funktion). Damit ein Element $s$ im nicht-assoziativen Raum $\mathbb{O}$ stabil bleibt, muss es symmetrisch bezüglich dieser Dualität liegen. Jede Asymmetrie würde eine Vorzugsrichtung erzeugen, die die Isomorphie der Untergitter bricht.
	
	Wir definieren das \textbf{Symmetrie-Funktional} $\Phi(s)$ als Differenz der Norm-Potentiale im Oktonionischen Raum:
	\begin{equation}
		\Phi(s) = \left| \|s\|_{\mathbb{O}}^2 - \|1-s\|_{\mathbb{O}}^2 \right|
	\end{equation}
	
	\begin{theorem}[Hurwitz-Stabilität]
		Damit die Divisionsalgebra $\mathbb{O}$ ihre Teilbarkeitseigenschaft behält, muss für alle generierenden Elemente $\rho$ gelten: $\Phi(\rho) = 0$.
	\end{theorem}
	
	\begin{proof}
		Sei $\rho = \sigma + i\gamma$. Wir setzen in $\Phi(\rho) = 0$ ein:
		\begin{align}
			|\rho|^2 &= |1-\rho|^2 \\
			\sigma^2 + \gamma^2 &= (1-\sigma)^2 + \gamma^2 \\
			\sigma^2 &= 1 - 2\sigma + \sigma^2 \\
			0 &= 1 - 2\sigma \quad \Rightarrow \quad \sigma = \frac{1}{2}
		\end{align}
	\end{proof}
	\noindent
	\textbf{Interpretation:} Die Riemann-Hypothese ist die Bedingung für das geometrische Gleichgewicht zwischen dem Nullpunkt (Leere) und der Eins (Einheit). Nur auf der kritischen Linie ist die \glqq Spannung\grqq{} im $E_8$-Gitter minimal.
	
	% ----------------------------------------------------------------------
	\section{Physikalisches Modell: Das modifizierte Planck-Gesetz}
	% ----------------------------------------------------------------------
	
	Wenn das Vakuum durch diese algebraische Struktur definiert ist, muss sich dies in der Zustandsdichte (DOS) zeigen. Wir folgen der Berry-Keating-Vermutung, dass die Nullstellen einem Spektrum $H=xp$ entsprechen.
	
	Wir postulieren ein modifiziertes Plancksches Strahlungsgesetz:
	\begin{equation}
		u_\nu^{\text{mod}}(\nu,T) = u_\nu^{0}(\nu,T) \cdot \Bigl[ 1 + \varepsilon \cdot M(z) \Bigr]
	\end{equation}
	mit $z = \ln(\nu/\nu_0)$. Die Modulationsfunktion $M(z)$ ist direkt aus der Gutzwiller-Spurformel abgeleitet:
	\begin{equation}
		M_{\text{RH}}(z) = \mathcal{N} \sum_{n=1}^{N} \frac{1}{\sqrt{\gamma_n^2 + 1/4}} \cos(\gamma_n z)
	\end{equation}
	Die Gewichtung $w_n \propto 1/|\rho_n|$ verhindert eine UV-Divergenz der Energie und spiegelt die Amplitudenabnahme im Hurwitz-Potential wider.
	
	% ----------------------------------------------------------------------
	\section{Methodik: Das Bamberg-Protokoll}
	% ----------------------------------------------------------------------
	
	Um diese Hypothese zu falsifizieren, definieren wir ein strenges Messprotokoll (\textit{Frozen Pipeline}), das p-Hacking ausschließt.
	
	\subsection{Phase I: Differenzmessung}
	Zur Eliminierung von Detektor-Artefakten (Responsivity $H(\nu)$) wird die Differenz zweier Schwarzkörper-Temperaturen $T_1, T_2$ gebildet:
	\begin{equation}
		S_{\text{diff}}(\nu) = \frac{S_{\text{mess}}(T_1)}{S_{\text{mess}}(T_2)} - \frac{u_\nu^0(T_1)}{u_\nu^0(T_2)}
	\end{equation}
	
	\subsection{Phase II: Statistische Inferenz (Gumbel)}
	Die extrahierten Residuen werden mittels Kreuzkorrelation gegen das Riemann-Template ($N=100.000$ Nullstellen, Datensatz \texttt{zeros6}) getestet.
	Die Teststatistik $T = \max_\tau |C(\tau)|$ folgt unter der Nullhypothese (reines Rauschen) einer \textbf{Gumbel-Verteilung}.
	Wir fordern eine Signifikanz von $p < 2.8 \times 10^{-7}$ ($5\sigma$) sowie erfolgreiche \textit{Adversarial Tests} gegen Zufallsmatrizen (GUE) und permutierte Nullstellen.
	
	% ----------------------------------------------------------------------
	\section{Fazit und Ausblick}
	% ----------------------------------------------------------------------
	
	Das hier vorgestellte Modell schlägt eine Brücke von der abstrakten Algebra des 19. Jahrhunderts zur experimentellen Physik des 21. Jahrhunderts.
	\begin{itemize}
		\item \textbf{Mathematisch:} Die Riemann-Hypothese ($\sigma=1/2$) ist notwendig, um die Symmetrie der Divisionsalgebren ($\mathbb{H} \to \mathbb{O}$) zu wahren.
		\item \textbf{Physikalisch:} Diese Symmetrie prägt die Feinstruktur des Vakuums und ist als Modulation der Hohlraumstrahlung messbar.
	\end{itemize}
	Wir rufen die Metrologie-Institute auf, existierende Rohdaten mittels des bereitgestellten Bamberg-Protokolls (Open Source) zu analysieren. Ein positiver Befund wäre der empirische Beweis, dass die Arithmetik der Primzahlen die fundamentale Grammatik der physikalischen Realität ist.
	
	% ----------------------------------------------------------------------
	\begin{thebibliography}{9}
		\small
		\bibitem{hurwitz1898} A. Hurwitz, \textit{Über die Composition der quadratischen Formen}, 1898.
		\bibitem{berry1999} M. V. Berry, J. P. Keating, \textit{The Riemann Zeros and Eigenvalue Asymptotics}, SIAM Review, 1999.
		\bibitem{baez2002} J. C. Baez, \textit{The Octonions}, Bull. Amer. Math. Soc., 2002.
		\bibitem{odlyzko} A. M. Odlyzko, \textit{The $10^{20}$-th zero of the Riemann zeta function}, 1989.
	\end{thebibliography}
	
\end{document}